% --- INICIO DE PLANTILLA PERSONALIZADA ---
\documentclass[oneside,11pt]{Latex/Classes/PhDthesisPSnPDF}

% --- Paquetes Originales ---
\usepackage{blindtext}
\usepackage{actuarialangle} 
\usepackage{tabularx}
\usepackage{float}
\usepackage{multirow, array}
\usepackage[usenames,dvipsnames,svgnames,table]{xcolor}
\usepackage{longtable}
\usepackage{latexsym,amsmath,amssymb,amsfonts}
\usepackage{enumitem}
\usepackage{xparse}
\usepackage{booktabs}
\usepackage[round, sort, numbers]{natbib}  
\usepackage{adjustbox}
\usepackage[utf8]{inputenc}
\usepackage{caption}
\usepackage{graphicx}
\usepackage{hyperref}
\usepackage{cleveref}

% Definición de colores
\definecolor{miblue}{RGB}{68, 114, 196}

% --- CAMBIO CRÍTICO 1: Usar \input para macros en el preámbulo ---
% \include da problemas antes del begin document. \input es seguro.
\input{Latex/Macros/MacroFile1}

% --- Definiciones de seguridad (por si la clase falla al cargar alguna variable) ---
\providecommand{\facultad}[1]{\def\lafacultad{#1}}
\providecommand{\escudofacultad}[1]{\def\elescudo{#1}}
\providecommand{\degree}[1]{\def\elgrado{#1}}
\providecommand{\director}[1]{\def\eldirector{#1}}
\providecommand{\lugar}[1]{\def\ellugar{#1}}
\providecommand{\degreedate}[1]{\def\lafecha{#1}}

% --- Configuración para bloques de código de R (Pandoc) ---

%%%%%%%%%%%%%%%%%%%%%%%%%%%%%%%%%%%%%%%%%%%%%%%%%%%%%%%%%%%%%%%%%%%%%%%%%%%%%%%%
%                         DATOS (Conectados a RMarkdown)                       %
%%%%%%%%%%%%%%%%%%%%%%%%%%%%%%%%%%%%%%%%%%%%%%%%%%%%%%%%%%%%%%%%%%%%%%%%%%%%%%%%
\title{Análisis de Juegos de Steam: ¿Cuáles son los mejores juegos?}
\author{Christopher Escalante \\ Gabriel Guevara \\ Noryen Harringhton}

% Lógica para Facultad: Si la defines en el YAML la usa, si no, usa la default

\facultad{Facultad de Ciencias Económicas y Sociales \\ Escuela de
Estadistica y Ciencias Actuariales}
  \escudofacultad{Latex/Classes/Escudos/faces}

\degree{Computación I}
\director{Jesus Ochoa \\ Oliver Triveño}
\degreedate{Abril 2026} 
\lugar{Caracas}

\portadatrue 

% Metadatos PDF
\keywords{}
\subject{}

% --- CORRECCIÓN PARA PANDOC NUEVO (Imágenes) ---
\makeatletter
\providecommand{\pandocbounded}[1]{#1}
\makeatother


%%%%%%%%%%%%%%%%%%%%%%%%%%%%%%%%%%%%%%%%%%%%%%%%%%%%%
%                   DOCUMENTO                       %
%%%%%%%%%%%%%%%%%%%%%%%%%%%%%%%%%%%%%%%%%%%%%%%%%%%%%
\begin{document}

\maketitle

%%%%%%%%%%%%%%%%%%%%%%%%%%%%%%%%%%%%%%%%%%%%%%%%%%%%%
%                  PRÓLOGO                          %
%%%%%%%%%%%%%%%%%%%%%%%%%%%%%%%%%%%%%%%%%%%%%%%%%%%%%
\frontmatter

% Puedes agregar dedicatorias desde el YAML si quieres, o dejarlas fijas aquí

%%%%%%%%%%%%%%%%%%%%%%%%%%%%%%%%%%%%%%%%%%%%%%%%%%%%%
%                   ÍNDICES                         %
%%%%%%%%%%%%%%%%%%%%%%%%%%%%%%%%%%%%%%%%%%%%%%%%%%%%%
\setcounter{secnumdepth}{3}
\setcounter{tocdepth}{3}

\tableofcontents



\cleardoublepage

  \cleardoublepage       % Empieza en página derecha (estándar en tesis)
  \chapter*{Resumen}     % Crea el título "Resumen" sin numeración (Capítulo *)
  \addcontentsline{toc}{chapter}{Resumen} % (Opcional) Lo añade al índice
  
  resumen             % Aquí se pega el texto que escribiste en el YAML
  
  \cleardoublepage

%%%%%%%%%%%%%%%%%%%%%%%%%%%%%%%%%%%%%%%%%%%%%%%%%%%%%
%                   CONTENIDO (El Corazón)          %
%%%%%%%%%%%%%%%%%%%%%%%%%%%%%%%%%%%%%%%%%%%%%%%%%%%%%
\mainmatter
\def\baselinestretch{1.5}

% --- CAMBIO CRÍTICO 2: Aquí RMarkdown inyecta todo tu texto ---
\chapter*{Introduccion}

La industria del entretenimiento interactivo ha experimentado una
transformación radical en las últimas dos décadas. Steam, la plataforma
de Valve Corporation, ha pasado de ser una simple tienda de
actualizaciones para \emph{Counter-Strike} a convertirse en el
ecosistema digital más grande del mundo para PC, albergando más de
50,000 títulos y alcanzando picos de 30 millones de usuarios
concurrentes diarios (Valve, 2024). Sin embargo, esta abundancia de
oferta ha generado un problema crítico: \textbf{la parálisis por
decisión y la competencia por la atención del usuario}.\\

Para un desarrollador independiente (Indie), lanzar un juego en Steam es
como arrojar una botella al océano. Para un gigante AAA, es una
inversión multimillonaria con riesgo incierto. Y en medio de ambos, los
juegos \emph{Free to Play} (F2P) han transformado las reglas del
mercado: ya no se compite solo por el precio de entrada, sino por el
\textbf{tiempo de atención} del jugador, un recurso escaso.\\

El presente trabajo tiene como objetivo descifrar los determinantes del
éxito en Steam. No se trata solo de ``qué juegos son los mejores'', sino
de \textbf{cómo compiten diferentes modelos de negocio} (Premium
vs.~F2P) y \textbf{escalas de producción} (Indie vs.~AAA) en un mercado
sobresaturado. Utilizando herramientas de análisis de datos en
\textbf{R} y técnicas estadísticas, descompondremos las variables que
predicen la valoración, la retención y el alcance de un videojuego.

\chapter{Planteamiento del Problema}

La plataforma Steam es el principal termómetro de la industria de
videojuegos para PC. Sin embargo, para los desarrolladores y
publicadores, el lanzamiento de un nuevo título conlleva un
\textbf{riesgo financiero significativo} debido a tres fenómenos
concurrentes:

\begin{enumerate}
\def\labelenumi{\arabic{enumi}.}
\item
  \textbf{Saturación del catálogo:} Con más de 50,000 juegos
  disponibles, la visibilidad orgánica es casi nula. La mayoría de los
  títulos nunca reciben suficientes reseñas para ser recomendados
  algorítmicamente.
\item
  \textbf{Guerra de escalas (Indie vs.~AAA):} Los grandes estudios (AAA)
  dominan el marketing y la producción gráfica, pero los juegos
  independientes (Indie) suelen innovar en mecánicas y construir
  comunidades más leales.
\item
  \textbf{Disrupción del Free to Play (F2P):} Títulos como \emph{Dota
  2}, \emph{Counter-Strike: Global Offensive} y \emph{Warframe} han
  demostrado que eliminar la barrera del precio inicial permite acumular
  decenas de millones de jugadores, acaparando cientos de horas de juego
  por usuario. Esto ha cambiado la expectativa del consumidor, que ahora
  compara un juego de \$60 no solo con otro de \$60, sino con 500 horas
  gratis en un F2P.
\end{enumerate}

Desde una perspectiva de ciencia de datos, el problema no es solo
``identificar juegos exitosos'', sino \textbf{modelar la incertidumbre}
y entender cómo interactúan las variables de escala (Indie/AAA), modelo
de negocio (Premium/F2P) y atributos del producto (género, precio,
soporte multiplataforma) para determinar el éxito.

\section{Preguntas de investigación}

Derivado de lo anterior, se plantean las siguientes preguntas de
investigación:

\textbf{1. Guerra de Escalas:} ¿Logran los juegos con la etiqueta
``Indie'' superar en ratio de reseñas positivas a los juegos de grandes
publicadores (AAA) a pesar de tener menor volumen de propietarios?

\textbf{2. Disrupción del F2P:} ¿Cómo se compara la retención real
(tiempo de juego) de los gigantes Free to Play frente a los títulos
Premium más exitosos? ¿El modelo gratuito genera ecosistemas donde solo
sobreviven los juegos como servicio?

\textbf{3. Economía del Jugador:} ¿Existe una correlación directa entre
el precio de un juego y su valoración, o los usuarios son más críticos
con los juegos más caros?

\textbf{4. Adquisición vs. Retención:} ¿Cuáles son los juegos con mayor
número de propietarios estimados? ¿Se traduce un alto número de
descargas en altos promedios de tiempo de juego?

\textbf{5. Barreras de Entrada:} ¿Qué porcentaje de incremento en ventas
estimadas y críticas positivas representa ofrecer un juego en
plataformas adicionales (Windows, Mac, Linux) o con soporte en múltiples
idiomas?

\textbf{6. Nichos Rentables:} Desde una óptica de desarrollo de
producto, ¿qué combinación de Género y Categoría presenta la menor
saturación en el mercado pero la mayor tasa de reseñas positivas?

\section{Justificación}

Este estudio es relevante por tres razones fundamentales:

\textbf{1.} Inteligencia de Mercado para Desarrolladores En un
ecosistema saturado, conocer qué atributos incrementan la probabilidad
de éxito permite diseñar productos más competitivos y minimizar riesgos
financieros.\\

\textbf{2.} Guía Estratégica para Jugadores y Consumidores El análisis
ayuda a identificar juegos con mejor relación calidad-precio, mayor
retención y mejores valoraciones, más allá de rankings superficiales.\\

\textbf{3.} Valor Académico y Profesional El proyecto integra:

\begin{itemize}
\item
  Ciencia de datos
\item
  Análisis estadístico
\item
  Visualización
\item
  RMarkdown
\item
  LaTeX
\item
  Trabajo colaborativo con GitHub
\end{itemize}

Esto lo convierte en un ejercicio formativo completo y aplicable al
ámbito profesional.

\section{Objetivos}

\subsection{Objetivo General}

Analizar los patrones de éxito de los videojuegos en Steam, evaluando
cómo variables como precio, género, tiempo de juego, modelo de negocio y
escala de producción influyen en la valoración y retención de los
usuarios

\subsection{Objetivos Específicos}

\begin{enumerate}
\item \textbf{Estacionalidad y Ciclo de Vida:} Describir la evolución temporal de las valoraciones y el tiempo de juego, identificando tendencias y puntos de inflexión en la serie de tiempo para determinar patrones estacionales favorables.

\item \textbf{Valoración y Retención por Categorías:} Determinar la influencia de las variables categóricas (Género, Plataforma, Escala de producción) en la distribución de las valoraciones y la retención de jugadores, identificando las mecánicas que generan mayor fidelidad.

\item \textbf{Análisis Competitivo (Indie vs. AAA):} Comparar el rendimiento en mercado de los juegos de grandes publicadores (AAA) frente a los desarrolladores independientes (Indie), evaluando si un menor presupuesto se traduce en menor calidad percibida o si los Indies lideran la satisfacción del usuario.

\item \textbf{Impacto del Modelo Free to Play (F2P):} Evaluar cómo los juegos gratuitos alteran la dinámica competitiva de la industria, analizando su capacidad para acaparar masivamente la cuota de mercado y el tiempo de juego, y cómo esta concentración afecta la viabilidad de los juegos de pago (Premium).

\item \textbf{Rentabilidad del Tiempo (Valor por Hora):} Analizar la relación entre el precio de venta y el tiempo mediano de juego para establecer un "costo por hora de entretenimiento" y entender la disposición a pagar del consumidor.

\item \textbf{Accesibilidad y Expansión de Mercado:} Evaluar cómo el soporte a múltiples plataformas (Windows, Mac, Linux) y la disponibilidad de idiomas afectan el tamaño del mercado potencial y las valoraciones de los usuarios.

\item \textbf{Reporte Técnico Reproducible:} Generar un reporte técnico en formato PDF que integre visualizaciones estadísticas y tablas resumen utilizando RMarkdown y LaTeX, garantizando la reproducibilidad del análisis.
\end{enumerate}

\section{Cobertura}

\subsection{Horizontal}

Describir la cobertura Horizontal del trabajo de investigación\\

\begin{table}[ht]
\centering
\caption{Variables consideradas} 
\resizebox{10cm}{!} {
\begin{tabular}{ccc}
  \hline
\ \ \ \ \ NOMBRE DE LA VARIABLE \\ 
  \hline
   VARIABLE 1 \\ 
VARIABLE 2\\ 
VARIABLE n\\ 
   \hline
\end{tabular}
}
\end{table}

\subsection{Vertical}

Describir la cobertura Vertical del trabajo de investigación

\section{Periodo de Referencia}

Describir la fuente y el Periodo de Referencia utilizado en el trabajo
de investigación

\section{Variables}

Describir la variable en estudio en el trabajo de investigación

\chapter{Marco Teórico}

\section{Bases Teóricas}

Resumen de las bases teoricas explicadas en el capitulo y su
importancia\\

\subsection{Definición 1}

Descripción de la definición 1\\

\subsection{Definición 2}

Descripción de la definición 2\\

\subsection{Definición n}

Descripción de la definición n\\

\chapter{Marco Metódico}

En esta sección se presentan las metodologías y/o test necesarios
relacionados con la práctica divulgada en el marco teórico, así como las
fuentes de datos.

\section{Bases metódicas utilizadas en el trabajo de investigación}

Desarrollo\\

\subsection{Procesamiento y Limpieza de Datos}

El proceso de Extracción, Transformación y Carga (ETL) consistió en:

\begin{enumerate}
\def\labelenumi{\arabic{enumi}.}
\item
  cargar datos desde el archivo 2.Steam.csv.
\item
  Limpieza: Eliminación de registros con valores nulos en variables
  críticas (price, positive\_ratings, average\_playtime).
\item
  Clasificación: Creación de variables derivadas:
\end{enumerate}

\begin{itemize}
\item
  tipo\_juego: ``Indie'' (si tiene etiqueta Indie o desarrollador indie
  conocido), ``AAA'' (si es de grandes estudios), ``Otros''.
\item
  modelo\_negocio: ``Free to Play'' si price == 0, ``Premium'' si price
  \textgreater{} 0.
\item
  ratio\_positivo: positive\_ratings / (positive\_ratings +
  negative\_ratings).
\item
  multiplataforma: ``Multiplataforma'' si platforms contiene punto y
  coma.
\end{itemize}

\begin{verbatim}
## [1] "Registros válidos: 27075 | Variables: 22"
\end{verbatim}

\subsection{Item 2}

Descripción\\

\subsubsection{Item 2.2}

Descripción\\

\subsection{Item n}

Descripción\\

\chapter{Análisis de Resultados}

\{\small Este capítulo presenta los resultados obtenidos al aplicar las
metodologías descritas en el capítulo anterior, así como del análisis
estadístico descriptivo de las variables estudiadas y las ventajas y
desventajas al aplicar dichos tests.\}

\section{Pregunta 1: Guerra de Escalas (Indie vs AAA}

\begin{center}\includegraphics[width=0.9\linewidth,]{trabajo-final_files/figure-latex/unnamed-chunk-1-1} \end{center}

Los juegos Indie muestran una mayor variabilidad en su ratio de reseñas
positivas, con valores que van desde muy bajos hasta muy altos.

Los AAA son más consistentes, con distribuciones más estrechas y
generalmente altas.

No se observa que los Indie superen claramente a los AAA, pero sí que
algunos Indie alcanzan niveles de valoración tan altos como los mejores
AAA.

\section{Presentación de resultados 2}

Descripción\\

\section{Presentación de resultados n}

Descripción\\

\chapter{Conclusiones y Recomendaciones}

\textbf{1.-} La comparación entre juegos AAA e Indie revela que los
títulos AAA presentan una distribución más concentrada y
consistentemente alta en el ratio de reseñas positivas, lo que sugiere
un estándar de calidad más homogéneo. En contraste, los juegos Indie
exhiben una mayor dispersión: aunque muchos presentan valoraciones
moderadas o bajas, existe un subconjunto significativo que alcanza
niveles de aceptación equivalentes o superiores a los AAA. Esto indica
que, si bien los Indie no superan a los AAA en términos centrales, sí
poseen un potencial de excelencia que desafía la hegemonía de los
grandes estudios.\\

\textbf{Recomendación} Dado que los juegos Indie presentan un potencial
de excelencia comparable al de los AAA, pero con una variabilidad mucho
mayor, se recomienda fomentar estrategias que reduzcan la dispersión
negativa ---como QA, marketing y diseño centrado en el usuario---
mientras se preserva la creatividad que caracteriza al sector
independiente.\\

\textbf{2.-} Conclusion 2 del trabajo de investigación.\\

\textbf{Recomendación} Recomendación 2 del trabajo de investigación.\\

\textbf{n.-} Conclusion n del trabajo de investigación.\\

\textbf{Recomendación} Recomendación n del trabajo de investigación.\\

\appendix
\renewcommand{\thechapter}{\Alph{chapter}}

\chapter{Anexo A}

\chapter{Anexo B}

\chapter{Anexo C}

%%%%%%%%%%%%%%%%%%%%%%%%%%%%%%%%%%%%%%%%%%%%%%%%%%%%%
%                   REFERENCIAS                     %
%%%%%%%%%%%%%%%%%%%%%%%%%%%%%%%%%%%%%%%%%%%%%%%%%%%%%
% Mantenemos la estructura natbib de tu plantilla
\renewcommand{\bibname}{Lista de Referencias}
\bibliographystyle{apalike} 
\bibliography{bibliografia.bib} 

%%%%%%%%%%%%%%%%%%%%%%%%%%%%%%%%%%%%%%%%%%%%%%%%%%%%%
%                   APÉNDICES                       %
%%%%%%%%%%%%%%%%%%%%%%%%%%%%%%%%%%%%%%%%%%%%%%%%%%%%%

\end{document}
